\documentclass[10pt,a4paper]{amsart}
\usepackage[margin=19mm]{geometry}
\usepackage{amsmath,amssymb,mathtools}
\usepackage[scr=rsfs,cal=euler]{mathalfa}
\usepackage{xfrac}
\usepackage[unicode,hidelinks,bookmarks=false]{hyperref}
\newcommand{\ie}{i.e.\ }
\newcommand{\cf}{cf.\ }
\newcommand{\resp}{respectively\ }
\newcommand{\Subsection}{\subsection}
\providecommand{\nobreakdash}{\nobreak\hbox{-}}
\setlength{\emergencystretch}{2em}
\multlinegap0pt
\usepackage[shortlabels]{enumitem}
\usepackage{mathtools}
\newtheorem{theorem}{Theorem}
\newtheorem{lemma}{Lemma}
\newtheorem{hypothesis}{Hypothesis}
\usepackage{cleveref}
\crefname{theorem}{Theorem}{Theorems}
\crefname{lemma}{Lemma}{Lemmas}
\crefname{hypothesis}{Hypothesis}{Hypothesiss}
\newcommand{\RSet}{\mathbb{R}}
\newcommand{\XSet}{\mathbb{X}}
\newcommand\blank{{\mkern 2mu\cdot\mkern 2mu}}
\title{A functional analytic theory for differential equations on Banach spaces with slowly evolving parameters}
\author{Dirk Doorakkers, Daniele Avitabile, Jan Bouwe van den Berg}
\date{Source: arXiv:2510.02893v2 (CC BY 4.0)}
\begin{document}
\maketitle

\noindent\emph{Source and licence.} A functional analytic theory for differential equations on Banach spaces with slowly evolving parameters. Version arXiv:2510.02893v2 (CC BY 4.0), distributed by arXiv under the Creative Commons Attribution 4.0 International licence, http://creativecommons.org/licenses/by/4.0/, as recorded in the arXiv licence metadata for this exact version and shown on the abstract page https://arxiv.org/abs/2510.02893v2. Full licence notice: \texttt{licenses/arxiv-2510.02893-CC-BY-4.0.txt}.\par\smallskip

\noindent\textit{Editorial scope.} Counted unit: the article's lemma lem:red-map-C1 (source0.tex lines 2417--2432). The article's complete original proof of it is delivered with it (source0.tex lines 2474--2555, one \texttt{proof} environment of 82 lines, which the article places away from the statement). No mathematics is written, repaired or re-derived: the statement and the proof are the article's own lines, in the article's own notation, and no retained line is substituted or re-flowed.  Retained uncounted as the source-local context the proof needs: lines 2108--2113; lines 2135--2149; lines 2221--2226; lines 2304--2316; lines 2374--2396.  Nothing was omitted from any retained span, so the package carries no omission record.  No retained line contains a citation, so the wrapper carries no bibliography and no bibliography entry.  No retained line contains a figure environment, an image-inclusion command or drawing code, so no image is delivered and the package declares no asset dependency.  Editorial formatting only, in addition: an AMS \texttt{amsart} preamble that restates the article's own declarations and macros used by the retained lines, namely the environments \texttt{theorem}, \texttt{lemma}, \texttt{hypothesis} on separate counters, together with the standard packages those lines need.  The retained environments print in this document as Theorem 1, Lemma 1, Hypothesis 1 in order of appearance, whereas the article prints the same items with the numbering of its own full document.  The complete native source of the article as distributed in the arXiv e-print for this exact version is bundled unchanged as \texttt{sources/186317/source0.tex}.
 \begin{equation}\label{ODE-system-red}
 \begin{aligned}
 \dot{x}(t) &= F(x(t),y(t)), \\
 \dot{y}(t) &= g(x(t),y(t)),
 \end{aligned}
 \end{equation}

 \begin{theorem}[Existence of Reduction Map]\label{thm:red-map-existence}
Assume System \cref{ODE-system-red} satisfies assumptions (S1)-(S2). Then there
exists a map $P \in C^0(\XSet \times \overline{V}, \overline{V})$ such that for each
$(\xi,\eta) \in \XSet \times \overline{V}$ it holds that 
\[
  \begin{aligned}
    & P(x(t; (\xi,\eta)),y(t;(\xi,\eta))) = y(t; (0,P(\xi,\eta))),  
      && \textit{for all $t \geq 0$,} \\
    & \lim_{t \rightarrow \infty} e^{\gamma t}  | y(t;(\xi,\eta)) - y(t; (0,P(\xi,\eta))) | = 0,
      && \textit{for any $0 \leq \gamma < \mu $,} \\
    & \sup \left\{ \frac{|\eta - P(\xi,\eta)|}{|\xi|} 
      \colon (\xi,\eta) \in \XSet \times \overline{V}, \xi \neq 0   \right\} < \infty.
  \end{aligned}
\]
 \end{theorem}

\begin{lemma}\label{lem:red-map-well-defined}
Under (S1), the mapping $\Lambda$ is well-defined on $\mathcal{E}^0$, and moreover
\[
  \| \Lambda(Q) \|_{\mathcal{E}} \leq \frac{K N_1}{\mu} (1+  \| Q \|_{\mathcal{E}}).
\]
\end{lemma}

\begin{lemma}\label{lem:red-map-props}
Assume (S1)-(S2) are valid. Suppose $Q \in  \mathcal{E}^0$ is such that $\Lambda(Q)=\Lambda$, and set $P(\xi,\eta) = \eta - Q(\xi,\eta)$. Then for every $(\xi,\eta) \in \XSet \times \overline{V}$ we have:
\begin{enumerate}
 \item For all $t \geq 0$ it holds that 
 \[
 \partial_t P(x(t; (\xi,\eta)),y(t;(\xi,\eta))) = g(0, P(x(t; (\xi,\eta)),y(t;(\xi,\eta))));
 \]
 \item $| Q(x(t; (\xi,\eta)),y(t;(\xi,\eta))) | \leq \frac{K^2 N_1}{\mu - K N_1} e^{-\mu t}  |\xi|$ for all $t \geq 0$, and thereby also
 \[
 \sup_{t \geq 0} e^{\mu t} | y(t;(\xi,\eta)) -  P(x(t; (\xi,\eta)),y(t;(\xi,\eta))) | \leq  \frac{K^2 N_1}{\mu - K N_1} |\xi|.
 \]
\end{enumerate}
\end{lemma}

 \begin{hypothesis}\label{hyp:red-map-l-derivatives}
Suppose $k \geq 2$, and that for some $1 \leq l \leq k$ we have:
\begin{enumerate}
\item 
  The functions $\partial^l_{(\xi,\eta)} x$ and $\partial^l_{(\xi,\eta) }y$
  exist and are continuous with respect to $(\xi,\eta)$. Further, there exist
  constants $K^x_l, K^y_l \geq 1$ such that for all $t \geq s \geq 0$ 
  \[
    \begin{aligned}
  & \| \partial^l_{(\xi,\eta) }x(t;(\xi,\eta)) \| 
    \leq K^x_l e^{-(\mu - l N_1) (t-s)} \|x(s;(\xi,\eta))\|, \\
  & \| \partial^l_{(\xi,\eta) }y(t;(\xi,\eta)) \| \leq 
  K^y_l e^{l N_1 (t-s)} \|y(s;(\xi,\eta))\|.
    \end{aligned}
  \]
\item 
  It holds $g \in C^k(U \times \overline{V}, \RSet^n)$. Further there exist 
  constants $N_1, \ldots, N_l$ such that $(l+1)N_1 < \mu$ and
\[
  \sup \{\|D^i g(x,y)\| : (x,y) \in U \times \overline{V} \} \leq N_i, \qquad \text{for all $ 1 \leq i \leq l$.}
\]
\end{enumerate}
\end{hypothesis}

 \begin{lemma}\label{lem:red-map-C1}
Suppose system \cref{ODE-system-red} satisfies the assumptions (S1)-(S2) from \cref{thm:red-map-existence}. Assume additionally that
\begin{enumerate}
\item \cref{hyp:red-map-l-derivatives} holds for $k=2$ and $l=1$
\item The map
 \begin{align*}
 D_y g : U \times \overline{V} &\rightarrow L(\RSet^n,\RSet^n) \, : \\
U \times \overline{V} \ni (x,y) &\mapsto D_y g(x, y) \in L( \RSet^n,\RSet^n)
 \end{align*}
is continuously differentiable at every point $(0,y) \in U \times V$ and there exists $N_2 > 0$ such that
\[
\sup_{x \in U, \, y_1,y_2 \in V} \| D_y g(x,y_2) - D_y g(0,y_1)  \| \leq N_2 (|x| + |y_2 - y_1|).
\]
\end{enumerate}
 Then for the map $P \in C^0(\XSet \times \overline{V}, \overline{V})$ from \cref{thm:red-map-existence} it holds moreover that $P \in  C^1(U \times \overline{V}, \overline{V})$.
 \end{lemma}

\begin{proof}[Proof of \cref{lem:red-map-C1}]
We firstly define a map
\[
Q^1 \in C^0(U \times \overline{V}, L(\XSet \times \RSet^n, \RSet^n)),
\]
which is now our candidate for the derivative of $(\xi,\eta) \mapsto \eta - P(\xi,\eta)$, as
\begin{align*}
Q^1(\xi,\eta) &:=  \int_0^{\infty} Z(0,s; (\xi,\eta)) \Bigg( D_y g(0,y(s;(\xi,\eta)) - q(s;(\xi,\eta))) \partial_{(\xi,\eta)}y(s; (\xi,\eta))  \\
&\qquad \qquad - D g(x(s;(\xi,\eta)),y(s;(\xi,\eta))) \begin{pmatrix} \partial_{(\xi,\eta)}x(s; (\xi,\eta)) \\ \partial_{(\xi,\eta)}y(s; (\xi,\eta)) \end{pmatrix} \Bigg) \, ds.
\end{align*}
We firstly have to show that $Q^1$ is well-defined in this manner. Using the
assumptions from \cref{lem:red-map-C1} we have the estimate
\begin{align*}
\| Q^1(\xi,\eta) \| &\leq  \int_0^{\infty} \| Z(0,s, (\xi,\eta)) \| \bigg( N_1 \| \partial_{(\xi,\eta)}x(s; (\xi,\eta)) \| \\
&\qquad \qquad \qquad + N_2 \left( |x(s; (\xi,\eta))| + |q(s; (\xi,\eta))| \right) \| \partial_{(\xi,\eta)}y(s; (\xi,\eta)) \| \bigg)  \, ds \\
&\leq \left( \frac{K^x_1 N_1}{\mu -  2 N_1} + \frac{K N_2 }{\mu - 2N_1} + \frac{K^2 K^y_1 N_1 N_2}{(\mu - K N_1)(\mu - 2 N_1)} \right) |\xi|,
\end{align*}
and thereby
\[
\sup_{\substack{(\xi,\eta) \in U \times \overline{V}; \\ \xi \neq 0}} \frac{\|Q^1(\xi,\eta)\|}{|\xi|} < \infty.
\]

It can be shown that $Q^1$ depends continuously on its arguments in a manner
comparable to the proof of \cref{lem:red-map-well-defined}. Therefore
\[
Q^1 \in \left\{ \tilde{Q} \in C^{0}(U \times \overline{V}, L(\XSet \times \RSet^n, \RSet^n)): \tilde{Q}(0, \blank) = 0 ;  \sup_{\substack{(\xi,\eta) \in U \times \overline{V}; \\ \xi \neq 0}} \frac{\|\tilde{Q}(\xi,\eta)\|}{|\xi|} < \infty  \right\}.
\]

Now for $q^1(t;\xi,\eta) := Q^1(x(t;(\xi,\eta)),y(t;(\xi,\eta)))$ set
\[
p^1(t;(\xi,\eta)) := \begin{pmatrix} 0 \\ I_{\RSet^n} \end{pmatrix} - q^1(t;(\xi,\eta)), \quad P^1(\xi,\eta) := p^1(0;(\xi,\eta)).
\]
It remains to show that $P^1$ is indeed the derivative of $P$. For this, we define for $t \geq 0$, $(\xi,\eta) \in U \times \overline{V}$ and $\Delta \in \XSet \times \RSet^n$:
\begin{align*}
a(t; (\xi,\eta), \Delta) :&= p(t; (\xi,\eta) + \Delta) - p(t; (\xi,\eta)) - p^1(t; (\xi,\eta)) \, \Delta.
\end{align*}
The quantity $a(t; (\eta,\xi), \Delta)$ has to satisfy the ODE
\begin{equation*}
\begin{aligned}
\partial_t \, a(t; (\xi,\eta), \Delta) &= g(0,p(t;(\xi,\eta)+\Delta)) - g(0,p(t;(\xi,\eta))) \\
&\qquad \qquad \qquad - D_y g(0,p(t;(\xi,\eta))) \, p^1(t;(\xi,\eta)) \, \Delta.
\end{aligned}
\end{equation*}

If we set
\begin{align*}
\hat{g}(t; (\xi,\eta), \Delta) :&= g(0,p(t;(\xi,\eta)+\Delta)) - g(0,p(t;(\xi,\eta))) \\
&\qquad \qquad \qquad - D_y g(0,p(t;(\xi,\eta))) \left( p(t;(\xi,\eta)+\Delta) - p(t;(\xi,\eta))  \right),
\end{align*}
we can rewrite the ODE for $a(t; (\xi,\eta), \Delta)$ as
\begin{equation}\label{eq:a-ODE}
\begin{aligned}
\partial_t \, a(t; (\xi,\eta), \Delta) &= D_y g(0,p(t;(\xi,\eta))) \, a(t; (\xi,\eta), \Delta) + \hat{g}(t; (\xi,\eta), \Delta) .
\end{aligned}
\end{equation}
Observe that by applying the Mean Value Theorem in integral form we may obtain
\[
| \hat{g}(t; (\xi,\eta), \Delta) | \leq \frac{N_2}{2} \left| p(t;(\xi,\eta)+\Delta) - p(t;(\xi,\eta))  \right|^2.
\]
By the first property \cref{lem:red-map-props}, and because of the second item of \cref{hyp:red-map-l-derivatives}, it holds that
\[
| p(t;(\xi,\eta)) | \leq | P(\xi,\eta) | e^{N_1|t|}.
\]
We therefore have for $\gamma > 2N_1$:
\[
\sup_{t \geq 0} e^{-\gamma t} | \hat{g}(t; (\xi,\eta), \Delta) | = o(|\Delta|).
\]

The relevant solutions to the ODE \cref{eq:a-ODE} satisfy the integral equation
\[
 a(t; (\xi,\eta), \Delta) =   \int_t^{\infty} Z(t,s; (\xi,\eta)) \, \hat{g}(s; (\xi,\eta), \Delta) \, ds,
\]
and therefore
\[
\sup_{t \geq 0} e^{-\gamma t} |  a(t; (\xi,\eta), \Delta)  | = o(|\Delta|).
\]
As $a(0; (\xi,\eta), \Delta) = P((\xi,\eta)+\Delta) - P(\xi,\eta) - P^1(\xi,\eta) \Delta$ we in particular have
\[
\left| P((\xi,\eta)+\Delta) - P(\xi,\eta) - P^1(\xi,\eta) \Delta \right| =   o(|\Delta|),
\]
and thereby $P^1 = DP$, as was desired. 
\end{proof}

\end{document}
